\section{二维随机变量函数的分布}

本节简单介绍二维随机变量的函数的分布。

本节要点：
\begin{itemize}
    \item 了解二维随机变量函数的概念；
    \item 了解几个独立同分布随机变量的叠加。
\end{itemize}

%============================================================
\subsection{二维随机变量函数的概念}

\begin{definition}[二维随机变量的函数]
设$\left( X,Y \right) $为二维随机变量，若$Z=g\left( X,Y \right) $依然为随机变量，则称$Z=g\left( X,Y \right) $为{\bf 二维随机变量$\left( X,Y \right) $的函数}。
\end{definition}

当$\left( X,Y \right) $为离散变量时，$Z=g\left( X,Y \right) $也为离散变量，其分布律可由列表的方式计算，注意需将相同值的概率进行合并。
当$\left( X,Y \right) $为连续变量时，$Z=g\left( X,Y \right) $可能是连续的、离散的，或者不定的。

%============================================================
\subsection{二维离散型随机变量的函数}

\begin{tcolorbox}
当$\left( X,Y \right) $为离散变量时，$Z=g\left( X,Y \right) $也为离散变量。
这种情况较为简单，重点是两个比较常见的分布的函数。
\end{tcolorbox}

\begin{theorem}
若随机变量$X,Y$独立且各自服从二项分布，则它们的和也服从二项分布：
\[
\begin{cases}
	X\sim B\left( m,p \right)\\
	Y\sim B\left( n,p \right)\\
\end{cases}\Rightarrow X+Y\sim B\left( m+n,p \right)
\]
\end{theorem}

\begin{theorem}
若随机变量$X,Y$独立且各自服从泊松分布，则它们的和也服从泊松分布：
\[
\begin{cases}
	X\sim P\left( \lambda _1 \right)\\
	Y\sim P\left( \lambda _2 \right)\\
\end{cases}\Rightarrow X+Y\sim P\left( \lambda _1+\lambda _2 \right)
\]
\end{theorem}

%============================================================
\subsection{二维连续型随机变量的函数}

\begin{theorem}[卷积公式]
设二维随机变量$\left( X,Y \right) $有密度函数$f=\left( x,y \right) $，若$Z=X+Y$，则$Z$的密度函数为：
\[
f_Z\left( Z \right) =\int_{-\infty}^{+\infty}{f\left( x,z-x \right) dx}=\int_{-\infty}^{+\infty}{f\left( z-y,y \right) dy}
\]
若$X,Y$相互独立，则：
\[
f_Z\left( Z \right) =\int_{-\infty}^{+\infty}{f\left( x \right) f\left( z-x \right) dx}=\int_{-\infty}^{+\infty}{f\left( y \right) f\left( z-y \right) dy}
\]
该公式称为{\bf 卷积公式}。
\end{theorem}

\begin{theorem}
设两个独立的随机变量$X,Y$各自服从正态分布，则它们的线性组合也服从正态分布：
\[
\begin{cases}
	X\sim N\left( \mu _1,{\sigma _1}^2 \right)\\
	Y\sim N\left( \mu _2,{\sigma _2}^2 \right)\\
\end{cases}\Rightarrow Z=aX+bY\sim N\left( a\mu _1+b\mu _2,a^2{\sigma _1}^2+b^2{\sigma _2}^2 \right)
\]
其中，$a,b$不能同时为0。
\end{theorem}

\begin{theorem}
设二维随机变量$\left( X,Y \right) $服从正态分布，则它们的线性组合也服从正态分布：
\begin{align*}
&\left( X,Y \right) \sim N\left( \mu _1,{\sigma _1}^2,\mu _2,{\sigma _2}^2,\rho \right) \\
&\Rightarrow \\
&Z=aX+bY\sim N\left( a\mu _1+b\mu _2,a^2{\sigma _1}^2+b^2{\sigma _2}^2 \right)
\end{align*}
其中，$a,b$不能同时为0。
\end{theorem}

~

关于二维随机变量函数更多的概念、定理和性质，见“教材\cite{book1}”。




